% 原创教学几何示意；四幅图在正文宽度内按两行两列排列，固定坐标与解析曲线。
\begin{tikzpicture}[x=1mm,y=1mm,font=\footnotesize]
  \foreach \xshift/\yshift/\title in {
    0/52/线性预测,55/52/非线性预测,
    0/0/线性分类,55/0/非线性分类}{
    \begin{scope}[xshift=\xshift mm,yshift=\yshift mm]
      \node[font=\heiti\small] at (15,37) {\title};
      \fill[BookPaper] (0,0) rectangle (30,28);
    \end{scope}
  }
  \foreach \xshift in {0,55}{
    \begin{scope}[xshift=\xshift mm,yshift=52mm]
      \draw[book arrow] (0,0)--(33,0) node[right] {$x$};
      \draw[book arrow] (0,0)--(0,30) node[above] {$\hat y$};
    \end{scope}
  }
  \foreach \xshift in {0,55}{
    \begin{scope}[xshift=\xshift mm]
      \draw[book arrow] (0,0)--(33,0) node[right] {$x_1$};
      \draw[book arrow] (0,0)--(0,30) node[above] {$x_2$};
    \end{scope}
  }
  \begin{scope}[yshift=52mm]
    \draw[FigureModel,very thick] (2,4)--(28,25);
  \end{scope}
  \begin{scope}[xshift=55mm,yshift=52mm]
    \draw[FigureModel,very thick]
      plot[domain=2:28,samples=60] ({\x},{4+0.032*(\x-2)^2});
  \end{scope}
  \begin{scope}
    \draw[FigureModel,thick] (2,4)--(28,24);
    \foreach \x/\y in {4/13,8/20,13/23,19/26}{
      \fill[FigureInput] (\x,\y) circle (1);
    }
    \foreach \x/\y in {9/4,17/7,24/10,28/17}{
      \fill[FigureLoss] (\x-0.9,\y-0.9) rectangle (\x+0.9,\y+0.9);
    }
  \end{scope}
  \begin{scope}[xshift=55mm]
    \draw[FigureModel,thick] (15,14) circle (8);
    \foreach \x/\y in {11/12,16/18,19/12,14/8}{
      \fill[FigureInput] (\x,\y) circle (1);
    }
    \foreach \x/\y in {4/6,5/25,16/26,28/22,27/5}{
      \fill[FigureLoss] (\x-0.9,\y-0.9) rectangle (\x+0.9,\y+0.9);
    }
  \end{scope}
  \node at (15,44) {固定增量};
  \node at (70,44) {增量随输入改变};
  \node at (15,-8) {直线分界};
  \node at (70,-8) {曲线分界};
\end{tikzpicture}
